\section{Literature Review and Research Status}
\label{sec:literature_review}

The intersection of artificial intelligence detection, factual verification, and computational morphology represents an active, rapidly evolving frontier in modern natural language processing. This section surveys current international and domestic research across these five foundational domains.

\subsection{AI-Generated Text Detection Methodologies}
Automated detection of machine-generated text has advanced through three successive methodological waves: metric-based statistical detectors, supervised neural classifiers, and perturbation-based curvature analyzers.

Early statistical techniques leveraged intrinsic token distribution properties, evaluating n-gram perplexity, rank order frequencies, and burstiness under assumed language model priors \cite{solaiman2019release}. These metrics capitalize on the empirical observation that autoregressive generation strategies (e.g., top-$p$ nucleus sampling, top-$k$ truncation) disproportionately select high-probability tokens from truncated distributions, yielding uniformly low perplexity trajectories compared to human text. Supervised neural classifiers subsequent to RoBERTa and DeBERTa approached detection as a standard binary sequence classification problem, learning discriminative stylistic artifacts directly from labeled datasets \cite{devlin2019bert, conneau2020xlmr}. However, supervised baselines suffer from severe distribution shifts: classifiers trained on older generator architectures (e.g., GPT-2) exhibit catastrophic performance degradation when tested on newer, instruction-tuned foundation models (e.g., LLaMA-3, Claude, GPT-4).

To eliminate the dependency on static training distributions, recent efforts have focused on zero-shot, training-free detection paradigms and comprehensive multi-generator benchmark robustness. Mitchell et al.\ \cite{mitchell2023detectgpt} introduced \textsc{DetectGPT}, positing that machine-generated texts occupy local optima within the model's log-probability curvature landscape. By applying mask-filling perturbations (via T5) and measuring the directional log-likelihood divergence between perturbed and original sequences, DetectGPT achieves competitive zero-shot detection. Hans et al.\ \cite{hans2024binoculars} developed \textsc{Binoculars}, which formulates detection as the normalized ratio of cross-perplexity between two distinct pre-trained foundation models, removing the latency of stochastic sampling. Meanwhile, Bao et al.\ \cite{bao2024fastdetectgpt} presented \textsc{Fast-DetectGPT}, utilizing analytical approximations of the Hessian curvature to accelerate inference speeds by three orders of magnitude. Expanding beyond single-model evaluations, Wang et al.\ \cite{wang2024raid} established \textsc{RAID}, a massive benchmark covering millions of generations across eleven contemporary LLMs, demonstrating that detector accuracy decays markedly under adversarial paraphrasing, prompt variations, and out-of-distribution generator shifts.

Despite their algorithmic elegance, virtually all contemporary detection benchmarks (e.g., HC3, M4, Ghostbuster \cite{verma2024ghostbuster}, RADAR \cite{hu2023radar}, RAID \cite{wang2024raid}) operate almost exclusively on English, German, or Mandarin Chinese. In low-resource agglutinative languages such as Kazakh, zero-shot curvature scoring fails because the surrogate scoring models (such as multilingual GPT-2 or mBERT) possess inadequate parametric vocabularies for Turkic languages, producing noisy, unstable log-probability surfaces that destroy the detection signal.

\subsection{Automated Fact Verification and Retrieval-Augmented Pipelines}
Automated fact verification emerged as a structured NLP discipline following the seminal release of the Fact Extraction and VERification (FEVER) benchmark by Thorne et al.\ \cite{thorne2018fever}. The standard FEVER pipeline decomposes claim verification into a three-stage architecture: (1) \textit{Document Retrieval}, identifying relevant Wikipedia articles from a large knowledge repository; (2) \textit{Sentence Evidence Extraction}, selecting the minimal set of sentences sufficient to substantiate or refute the claim; and (3) \textit{Claim Verification / Natural Language Inference}, predicting a three-way veracity label: \textsc{Supported}, \textsc{Refuted}, or \textsc{Not Enough Info} (NEI).

Subsequent benchmarking expanded the FEVER paradigm into specialized domains and real-world multi-source contexts. FEVER 2.0 \cite{thorne2019fever2} introduced adversarial human attacks; MultiFC \cite{augenstein2019multifc} and Snopes-based datasets gathered claims from public fact-checking organizations; and CLIMATE-FEVER \cite{diggelmann2020climatefever} explored climate change assertions characterized by subtle scientific nuances and conflicting evidence. In the retrieval tier, early systems relied on sparse inverted index mechanisms, notably BM25 \cite{robertson2009bm25}. The introduction of Dense Passage Retrieval (DPR) \cite{karpukhin2020dpr} and unsupervised contrastive retrievers (Contriever \cite{izacard2022mcontriever}) established dense neural embeddings as the dominant standard for semantic matching. More recently, multi-granularity representations such as \textsc{M3-Embedding} \cite{chen-etal-2024-m3} have unified dense, lexical, and multi-vector retrieval across more than 100 languages via self-knowledge distillation. Furthermore, Zheng et al.\ \cite{zheng-etal-2024-evidence} demonstrated that evidence retrieval accuracy forms the primary empirical bottleneck in verification systems, establishing that high-precision retrieval is decisive for downstream classification fidelity.

Concurrently, the rapid evolution of foundation models has catalyzed specialized architectures for LLM fact verification and hallucination auditing. Tang et al.\ \cite{tang-etal-2024-minicheck} developed \textsc{MiniCheck}, synthesizing lightweight, high-performance verifiers trained explicitly on grounding documents to achieve reliable fact-checking at a fraction of full LLM inference overhead. Guan et al.\ \cite{guan-etal-2024-language} systematically analyzed hallucination phenomena, showing that while LLMs readily generate spurious statements during open production, their internal attention patterns can excel at discriminating factual inconsistencies when supplied with explicit context. To resolve complex or ambiguous claims, Yue et al.\ \cite{yue-etal-2024-retrieval} proposed retrieval-augmented fact verification through contrastive argument synthesis, generating counter-arguments to expose unsupported assertions. Meanwhile, fine-grained evaluation benchmarks such as \textsc{Factcheck-Bench} \cite{wang-etal-2024-factcheck-bench} and \textsc{FactLens} \cite{mitra-etal-2025-factlens} have introduced diagnostic frameworks to evaluate atomic claim decomposition, reasoning shortcuts, and multi-step evidence integration.

However, non-English fact verification benchmarks remain exceptionally sparse. While multilingual extensions such as X-FEVER, Multilingual FEVER, and Factcheck-Bench \cite{wang-etal-2024-factcheck-bench} incorporate translation-based evaluations, they overwhelmingly rely on synthetic machine translation of English benchmarks. These translated datasets systematically obscure indigenous cultural phenomena, local legal nuances, and native syntactic structures. Most critically, until this dissertation, there has been no dedicated, human-annotated fact verification benchmark for Kazakh or any other Central Asian language.

\subsection{Computational Morphology for Agglutinative Languages}
Computational modeling of agglutinative morphology has historically relied on two distinct traditions: formal symbolic rule engines and statistical/neural subword segmenters.

In the symbolic paradigm, Finite-State Transducers (FSTs) based on the two-level morphology framework (Koskenniemi) provide exact, bidirectional mappings between surface wordforms and grammatical analyses \cite{tyers2012free, khoshhal2020kazakh}. The Apertium collaborative platform developed open-source morphological analyzers for Kazakh, utilizing extensive dictionaries and morphotactic continuation lexicons to handle vowel harmony (palatal/labial harmony) and consonant assimilation \cite{tyers2012free}. While symbolic FSTs achieve near-perfect precision on in-vocabulary lexical stems, their integration with modern deep learning representations has been severely constrained by non-differentiable execution boundaries and brittleness when encountering out-of-vocabulary named entities or internet slang.

In the statistical paradigm, unsupervised subword algorithms such as Byte-Pair Encoding (BPE), WordPiece, and SentencePiece dominate modern language modeling. Although computationally efficient, recent studies have demonstrated that subwords introduce severe cognitive biases in agglutinative contexts. Khurana et al.\ \cite{khurana2024subwords} demonstrated that subword tokenizers frequently partition grammatical suffixes in Turkic and Uralic languages across arbitrary phoneme boundaries, completely obscuring syntactic dependencies and negation allomorphs. Hybrid frameworks that explicitly inject symbolic morphological features into neural embedding layers represent an urgent, under-explored research trajectory.

In parallel, the broader NLP landscape for Kazakh and Central Asian languages has witnessed vital foundational contributions. The \textsc{Belebele} benchmark \cite{bandarkar-etal-2024-belebele} established a massively parallel reading comprehension suite across 122 language variants including Kazakh. Togmanov et al.\ \cite{togmanov-etal-2025-kazmmlu} formulated \textsc{KazMMLU}, providing the first comprehensive evaluation of foundation models on Kazakh, Russian, and regional knowledge of Kazakhstan across 56 subjects, highlighting dramatic performance gaps between English and regional languages. Laiyk et al.\ \cite{laiyk-etal-2025-instruction} demonstrated the efficacy of instruction tuning on curated public government and cultural corpora in low-resource Kazakh. In specialized tasks and pre-training, Yeshpanov and Varol \cite{yeshpanov2024kazsandra} constructed \textsc{KazSAnDRA} for fine-grained sentiment analysis, Sagyndyk et al.\ \cite{sagyndyk2025kazroberta} introduced \textsc{Kaz-RoBERTa} for contextual conversational understanding, and Tukenov \cite{tukenov2026sozkz} released \textsc{SozKZ}, training efficient small language models natively from scratch. Despite these strides, existing resources focus predominantly on downstream classification or general pre-training; none systematically model the interaction between complex agglutinative suffix chains and multi-stage factual verification.

\subsection{Shortcut Learning and the ``Hard NEI'' Pathology}
A pervasive failure mode in contemporary natural language inference and fact-checking models is their vulnerability to \textit{shortcut learning} (also known as spurious correlation exploitation) \cite{wang-etal-2024-factcheck-bench}. In standard FEVER benchmarks, instances belonging to the \textsc{Not Enough Info} (NEI) class are typically created by pairing claims with unrelated passages retrieved via random sampling. Consequently, neural verifiers rapidly discover a superficial statistical shortcut: if a claim and an evidence passage exhibit low surface token overlap or distant semantic embeddings, the pair can be reliably classified as \textsc{NEI} with high training accuracy.

Recent diagnostic evaluations demonstrate that automatic fact-checking models exploit shallow heuristics---such as sentiment valence and superficial named entity matches---rather than genuine reasoning chains \cite{wang-etal-2024-factcheck-bench}. To enforce causal fidelity, Si et al.\ \cite{si-etal-2024-checkwhy} introduced \textsc{CheckWhy}, structuring fact verification around causal argument graphs that require models to distinguish valid causal links from spurious surface associations. In real-world media environments, false or unverified claims are rarely disjoint from authentic facts: malicious disinformation and subtle rumors routinely retain high lexical overlap with genuine news reports, altering only temporal anchors, quantitative bounds, evidential modalities, or logical scopes. When evaluated on such challenging instances, standard cross-encoders suffer catastrophic performance collapse, predicting \textsc{Supported} merely because the candidate evidence shares extensive topical vocabulary with the claim. Neutralizing this pathology requires explicit counterfactual curation protocols that demand fine-grained morphological reasoning.

\subsection{Explainable AI and Multidimensional Trust Frameworks}
In high-stakes domains such as investigative journalism, intelligence auditing, and content moderation, automated black-box predictions are fundamentally insufficient. Human auditors require transparent, interpretable rationales and calibrated uncertainty metrics before accepting model recommendations.

Feature attribution methods such as Integrated Gradients, Grad-CAM \cite{selvaraju2017gradcam}, and SHAP \cite{lundberg2017shap} provide post-hoc saliency heatmaps highlighting influential tokens. In retrieval-grounded systems, recent research demonstrates that hybrid retrieval combined with explicit evidence highlighting provides human auditors with bounded, verifiable context; for instance, Mudunuri et al.\ \cite{mudunuri-etal-2026-bounded} showed that bounding conversational memory with hybrid retrieval and evidence span highlighting reduces drift and elevates user trust during interactive verification. However, in low-resource agglutinative languages, token-level heatmaps are notoriously difficult to interpret when subwords segment suffixes into non-intuitive byte fragments. Furthermore, existing NLP benchmarks evaluate veracity (true vs.\ false) and generation origin (human vs.\ AI) in complete isolation. As generative AI becomes widely adopted by legitimate authors and malicious actors alike, automated decision systems require a unified, multi-criteria trust space that simultaneously captures factual accuracy and generative provenance, establishing an interpretable foundation for human-in-the-loop verification \cite{khashabi2021genie}.
